\section{Laplacians and tree packings}\label{sec:preliminaries}

Unless a base is specified, $\log$ denotes the natural logarithm.

All graphs in the proof are finite undirected loopless multigraphs.
They have a common vertex set when one is said to be a subgraph of
another. Each parallel edge has its own index.
For a partition \(\mathcal P\) of the vertex set, let
\(E_H(\mathcal P)\) denote the edges whose endpoints lie in different
parts. A graph is \(r\)-tree-packable if it contains \(r\) pairwise
edge-disjoint spanning trees. On a singleton vertex set this property
is taken to hold for every positive integer \(r\).

\subsection{Matrix conventions}

Choose an orientation of each edge solely to define
\(b_e=\mathbf1_u-\mathbf1_v\in\R^V\), and put
\[
 L_H=\sum_{e\in E(H)}b_eb_e^{\mathsf T}.
\]
Changing an orientation changes no matrix or resistance used below.
For a real symmetric positive definite matrix \(X\), define
\[
 R_X(e)=b_e^{\mathsf T}X^{-1}b_e.
\]
The order \(A\preceq B\) means that \(B-A\) is positive semidefinite.
We use the operator norm for matrices unless otherwise specified.
For every nonempty proper \(S\subset V\),
\[
 \mathbf1_S^{\mathsf T}L_H\mathbf1_S=|\delta_H(S)|.
\]
An inequality required only for these indicator vectors is called a
\emph{cut inequality}. It need not be a semidefinite inequality.
In particular, no cut inequality is imposed at \(S=V\).

\begin{lemma}\label{lem:resistance-order}
For \(X\succ0\), a column \(b\), and \(t>0\),
\[
 b^{\mathsf T}X^{-1}b\le t
 \quad\Longleftrightarrow\quad
 bb^{\mathsf T}\preceq tX.
\]
If \(X\succeq Y\succ0\), then
\(b^{\mathsf T}X^{-1}b\le b^{\mathsf T}Y^{-1}b\).
\end{lemma}
\begin{proof}
Congruence by \(X^{-1/2}\) reduces the first assertion to the fact
that the only possible nonzero eigenvalue of \(aa^{\mathsf T}\) is
\(\|a\|^2\), with \(a=X^{-1/2}b\).
For the second, \(Y^{-1/2}XY^{-1/2}\succeq I\); diagonalization and
inversion give \(X^{-1}\preceq Y^{-1}\).
\end{proof}

When a Laplacian is inverted, we will explicitly restrict it to
\[
 \mathbf1^\perp=\{x\in\R^V:\textstyle\sum_{v\in V}x_v=0\}.
\]
The shortcut matrices \(X\) used in extraction are positive definite
on the full space \(\R^V\), so their inverses require no such convention.

\begin{lemma}[Incidence rank]\label{lem:incidence-rank}
If \(F\) has \(n\) vertices and \(c\) connected components, its
incidence columns span a space of dimension \(n-c\).
In particular the columns of a spanning tree form a basis of
\(\mathbf1^\perp\).
\end{lemma}
\begin{proof}
A vector is orthogonal to every incidence column precisely when
its coordinates agree at the endpoints of each edge, equivalently
when it is constant on every connected component.
The orthogonal complement thus has dimension \(c\).
A tree on \(n\) vertices has \(n-1\) edges, as follows by deleting
leaves successively, so its spanning columns form a basis.
\end{proof}

\subsection{The partition criterion}

\begin{theorem}[Nash--Williams--Tutte]\label{thm:packing}
For a positive integer \(r\), a graph \(H\) contains \(r\)
edge-disjoint spanning trees if and only if
\[
 |E_H(\mathcal P)|\ge r(|\mathcal P|-1)
 \qquad\text{for every partition }\mathcal P\text{ of }V(H).
\]
\end{theorem}

We include an exchange proof of this classical
criterion~\cite{NashWilliams,Tutte}.
The insertion and cycle exchanges below are the graphic form of
the circuit exchanges in Edmonds's matroid-partition
proof~\cite[Section~1.6]{Edmonds1965}.

\begin{proof}
The necessity follows by counting the crossings of each tree.
For sufficiency, choose \(r\) disjoint forests whose total number
\(m\) of edges is maximum. A \emph{state} is such a choice.
Starting from this state, allow a move that adds an unused edge to
one forest and removes an edge of the resulting cycle.

In every reachable state, every unused edge closes a cycle with
each forest: otherwise adding it would increase \(m\).
Let \(U\) be the set of edges that are unused in at least one reachable
state. If an unused edge is added to a forest, every edge of the
resulting cycle lies in \(U\), since any one of its original edges
can be made unused by the corresponding move.
Such a move preserves the components of that forest restricted to
its edges in \(U\). It also leaves the other forests unchanged.

Fix \(e\in U\). In a state where \(e\) is unused, its endpoints are
connected by edges in \(U\) in each forest. The component invariance
just proved carries this statement to the initial state.
Consequently each component of \((V(H),U)\) is spanned by the
\(U\)-edges of every forest.
If there are \(d\) such components, this accounts for \(r(n-d)\)
used edges, where \(n=|V(H)|\).
Every edge between these components is used, since every unused
edge belongs to \(U\). Their number is therefore at most
\[
 m-r(n-d).
\]
If \(m<r(n-1)\), this is smaller than \(r(d-1)\), contrary to the
partition hypothesis. Thus \(m=r(n-1)\); each of the \(r\) forests
has \(n-1\) edges and is a spanning tree.
\end{proof}

\begin{lemma}[Union and quotient properties]\label{lem:packing-parts}
Fix a positive integer \(q\) and a graph \(F\).
\begin{enumerate}[label=\textup{(\roman*)}]
\item If \(F[A]\) and \(F[B]\) are \(q\)-tree-packable and
\(A\cap B\ne\varnothing\), then \(F[A\cup B]\) is
\(q\)-tree-packable.
\item The inclusion-maximal \(q\)-tree-packable vertex sets form a
partition \(\mathcal P(F)\). If \(F\subseteq F'\), then
\(\mathcal P(F')\) coarsens \(\mathcal P(F)\).
\item If \(|\mathcal P(F)|>1\), then
\[
 |E_F(\mathcal P(F))|<q(|\mathcal P(F)|-1).
\]
\end{enumerate}
\end{lemma}
\begin{proof}
For (i), enumerate the \(q\) trees on each of \(A\) and \(B\).
In the \(i\)-th tree on \(B\), contract all vertices of \(A\cap B\)
to a single vertex, discard loops, and choose a spanning tree of
the resulting connected multigraph. Its selected edges, together
with the \(i\)-th tree on \(A\), form a tree on \(A\cup B\).
Every selected edge has an endpoint outside \(A\), so no such edge
lies in any tree on \(A\). The selections in different colors
remain disjoint. This proves (i).

Every singleton is packable, and finiteness gives maximal sets.
Part (i) shows that two maximal sets cannot intersect unless equal.
Every previously packable set remains packable when edges are added,
which proves (ii).

For (iii), contract the parts and remove loops.
If its crossing count were at least \(q(d-1)\), where
\(d=|\mathcal P(F)|>1\), the quotient would have a set
\(W\) of at least two vertices with
\(|E(W)|\ge q(|W|-1)\).
Choose such \(W\) inclusion-minimal. For each proper nonempty
\(J\subset W\),
\[
 |E(J)|\le q(|J|-1);
\]
this also holds for a singleton because loops were removed.
For a partition \(W=J_1\sqcup\cdots\sqcup J_t\), with \(t\ge2\),
subtraction gives at least
\[
 q(|W|-1)-\sum_{\ell=1}^t q(|J_\ell|-1)=q(t-1)
\]
edges between its parts. Theorem~\ref{thm:packing} supplies \(q\)
quotient trees on \(W\).
Attach to each color a packed tree inside each of its quotient
vertices. This gives \(q\) disjoint spanning trees on their union
in \(F\), contradicting maximality of the parts.
\end{proof}
